%!TEX root = ../KTheorie_SS15/K-Theorie.tex
\RequirePackage[thinlines]{easybmat} %-- muss aufgrund eines Bugs vor etex und tikz geladen werden
\documentclass[a4paper, twoside, headsepline, index=totoc,toc=listof, fontsize=10pt, cleardoublepage=empty, headinclude, DIV=12, BCOR=5mm, titlepage,draft]{scrartcl}
\usepackage{scrtime} % KOMA, Uhrzeit ermoeglicht
\usepackage{scrpage2} % wie fancyhdr, nur optimiert auf KOMA-Skript, leicht andere Syntax
\usepackage{etoolbox}
\usepackage{letltxmacro}
\usepackage{ifthen}

%--Farbdefinitionen
%-- muss vor tikz geladen werden
\usepackage[usenames, table, x11names]{xcolor}
\definecolor{dark_gray}{gray}{0.45}
\definecolor{light_gray}{gray}{0.6}
\usepackage[final]{graphicx}

%--Zum Zeichnen
%-- muss vor polyglossia bzw. babel geladen werden
\usepackage{tikz}
\usepackage{tikz-cd}
\usetikzlibrary{external}
\tikzset{>=latex}
\usetikzlibrary{shapes,arrows,intersections}
\usetikzlibrary{calc,3d}
\usetikzlibrary{decorations.pathreplacing,decorations.markings, decorations.pathmorphing}
\usetikzlibrary{angles}

%-- Konfiguration von tikzexternalize
\tikzexternalize[prefix=tikz/,up to date check=diff]
\pgfkeys{/pgf/images/include external/.code=\includegraphics{#1}}
\tikzset{external/system call={lualatex \tikzexternalcheckshellescape -halt-on-error -interaction=batchmode --shell-escape -jobname "\image" "\texsource"}}


%-- tikzexternalize fuer tikzcd deaktivieren, da inkompatibel
\AtBeginEnvironment{tikzcd}{\tikzexternaldisable}
\AtEndEnvironment{tikzcd}{\tikzexternalenable}

%-- um Inkompatibilitaeten von quotes und polyglossia bzw. babel zu vermeiden
\tikzset{
  every picture/.append style={
    execute at begin picture={\shorthandoff{"}},
    execute at end picture={\shorthandon{"}}
  }
}
\usetikzlibrary{quotes}
\usepackage{pgfplots}
\usepgfplotslibrary{colormaps}
\pgfplotsset{compat=1.12}



%-- Mathesymbole etc.
\usepackage{mathtools} % beinhaltet amsmath
\mathtoolsset{showonlyrefs, centercolon}
\newtagform{brackets}[\textbf]{[}{]}
\usetagform{brackets}
\usepackage{wasysym} % zusätzliche Symbole
\usepackage{amssymb} % zusätzliche Symbole
% \usepackage{latexsym} % zusätzliche Symbole
\usepackage{stmaryrd} % für Blitz
\usepackage{nicefrac} % schräge Brüche
\usepackage{faktor}
\newcommand{\Faktor}[1]{\faktor[\textstyle]{#1}}
\usepackage{cancel} % Befehle zum Durchstreichen
\usepackage{mathdots} % Verbesserung von Punkten wie zB \ldots
\usepackage{mathrsfs} % das X wurde benötigt

%-- einzelne Symbole aus dem mathx-package
\DeclareFontFamily{U}{mathx}{\hyphenchar\font45}
\DeclareFontShape{U}{mathx}{m}{n}{<-> mathx10}{}
\DeclareSymbolFont{mathx}{U}{mathx}{m}{n}
\DeclareMathSymbol{\bigplus}{1}{mathx}{"90}

%-- 
\DeclareSymbolFont{bbold}{U}{bbold}{m}{n}
\DeclareSymbolFontAlphabet{\mathbbold}{bbold}
\newcommand{\mathds}[1]{\mathbb{#1}} % Um Kompatibilitaet mit frueheren Benutzung von dsfont herzustellen bzw zusätzliche Kapselung
\newcommand{\ind}{\mathbbold{1}} % charakteristische-Funktion-Eins

\def\mathul#1#2{\color{#1}\underline{{\color{black}#2}}\color{black}} %farbiges Untersteichen im Mathe-Modus
\renewcommand{\le}{\leqslant}
\renewcommand{\ge}{\geqslant}

%-- Underbrace als Befehl in LaTeX-Syntax (und ohne Spacingprobleme mit nachfolgenden Operatoren...)
\newcommand{\Underbrace}[2]{{\underbrace{#1}_{#2}}}

%-- Alles was mit Schrift und XeTeX zu tun hat
\usepackage[euler-digits]{eulervm}
\usepackage[no-math]{fontspec}
\usepackage{polyglossia} % moderner babel-ersatz
\setmainlanguage[spelling=new,babelshorthands=true]{german}
\shorthandoff{"}
\setotherlanguage{english}
% \newcommand\glqq{"}
% \newcommand\grqq{"}
\defaultfontfeatures{Mapping=tex-text, WordSpace={1.2}, Ligatures={Required,Common,Contextual}} %

%-- Alle hier aufgeführten Schriftarten sind in jeder TeXlive-Installation vorhanden!!
%-- Hauptschriftart festlegen
\setmainfont{texgyrepagella}[Extension=.otf,UprightFont=*-regular,BoldFont=*-bold,ItalicFont=*-italic,BoldItalicFont=*-bolditalic,ItalicFeatures={Style=Historic}]
% \setmainfont[BoldFont={* Bold}, ItalicFont={* Light Italic},SmallCapsFont={Linux Libertine O}, SmallCapsFeatures={Letters=SmallCaps}]{Source Sans Pro}


%-- Serifenlose Schriftart festlegen
% \setsansfont{Roboto}[Scale=MatchUppercase,Extension=.ttf,UprightFont=*-Regular,BoldFont=*-Medium,ItalicFont=*-RegularItalic,BoldItalicFont=*-MediumItalic]
% \setsansfont{RobotoSlab}[Scale=MatchUppercase,Extension=.ttf,UprightFont=*-Regular,BoldFont=*-Bold,ItalicFont=*-Thin,SmallCapsFont=*-Regular]
% \setsansfont{FiraSans}[Extension=.otf, UprightFont=*-Regular, BoldFont=*-Medium, ItalicFont=*-RegularItalic, BoldItalicFont=*-MediumItalic]
% \setsansfont{SourceSansPro}[Scale=MatchUppercase,Extension=.otf,UprightFont=*-Regular,BoldFont=*-Semibold,ItalicFont=*-LightIt,BoldItalicFont=*-SemiboldIt]
% \setsansfont{texgyreadventor}[Scale=MatchUppercase, Extension=.otf, UprightFont=*-regular, BoldFont=*-bold, ItalicFont=*-italic,BoldItalicFont=*-bolditalic]
% \setsansfont{Kurier}[Scale=MatchUppercase, Extension=.otf, UprightFont=*Medium-Regular, BoldFont=*Heavy-Regular, ItalicFont=*Medium-Italic,BoldItalicFont=*Heavy-Italic]
% \setsansfont{Iwona}[Scale=MatchUppercase, Extension=.otf, UprightFont=*Medium-Regular, BoldFont=*Heavy-Regular, ItalicFont=*Medium-Italic,BoldItalicFont=*Heavy-Italic]
\setsansfont{LinBiolinum}[Scale=MatchUppercase, Extension=.otf, UprightFont=*_R, BoldFont=*_RB, ItalicFont=*_RI,BoldItalicFont=*_RBO]


% \newfontfamily{\geometric}{texgyreheros}[Scale=MatchLowercase, Extension=.otf, UprightFont=*-regular, BoldFont=*-bold, ItalicFont=*-italic]
% \addtokomafont{disposition}{\geometric}
\setmonofont{SourceCodePro}[Scale=0.9,Extension=.otf,UprightFont=*-Regular, BoldFont=*-Semibold, ItalicFont=*-Light]
\usepackage{xltxtra}
\usepackage{fontawesome}
\usepackage[final]{microtype}


%--Mixed
\usepackage[shortlabels,inline]{enumitem}
\setlist[itemize,1]{label=\faCaretRight}
\setlist[enumerate]{font=\bfseries}
\usepackage[german=quotes]{csquotes}
\usepackage{makeidx}
\usepackage{booktabs}
\usepackage{wrapfig}
\usepackage{float}
\usepackage[margin=10pt, font=small, labelfont={sf, bf}, format=plain, indention=1em]{caption}
\captionsetup[wrapfigure]{name=Abb. }
\usepackage{stackrel}
\usepackage{multicol}

\flushbottom


%--Unterstreichung
\usepackage[normalem]{ulem}
\setlength{\ULdepth}{1.8pt}

%--Indexverarbeitung
\newcommand{\bet}[1]{\emph{#1}}
\newcommand{\Index}[1]{\bet{#1}\index{#1}}
\makeindex
\setindexpreamble{{\noindent \itshape Die \emph{Seitenzahlen} sind mit Hyperlinks zu den entsprechenden Seiten versehen, also anklickbar} \faHandUp \par \bigskip}
\renewcommand{\indexpagestyle}{scrheadings}


%--Marginnote & todonotes
\deffootnote[1.5em]{1.5em}{1.5em}{\textsuperscript{\thefootnotemark}\ }
\usepackage[fulladjust]{marginnote}
\renewcommand*{\marginfont}{\color{dark_gray} \itshape \footnotesize}
\usepackage[textsize=small]{todonotes}
\usepackage{ragged2e}
\renewcommand*{\raggedleftmarginnote}{\RaggedLeft}
\renewcommand*{\raggedrightmarginnote}{\RaggedRight}

%--Todonotes mit tikz/externalize kompatibel machen
\LetLtxMacro{\oldtodo}{\todo}
\renewcommand{\todo}[2][]{\tikzexternaldisable\oldtodo[#1]{#2}\tikzexternalenable}
\LetLtxMacro{\oldmissingfigure}{\missingfigure}
\renewcommand{\missingfigure}[2][]{\tikzexternaldisable\oldmissingfigure[{#1}]{#2}\tikzexternalenable}

%--Konfiguration von Hyperref pdfstartview=FitH, 
\usepackage[hidelinks, pdfpagelabels,  bookmarksopen=true, bookmarksnumbered=true, linkcolor=black, urlcolor=SkyBlue2, plainpages=false,pagebackref, citecolor=black, hypertexnames=true, pdfauthor={Jannes Bantje}, pdfborderstyle={/S/U}, linkbordercolor=SkyBlue2, colorlinks=false,final]{hyperref}


\newcommand{\appendLink}[1]{#1\,\faExternalLink}
\newcommand{\hrefsym}[2]{\href{#1}{\texttt{\appendLink{#2}}}}
\newcommand{\hrefsymmail}[2]{\href{#1}{\texttt{\faEnvelope\,#2}}}
\renewcommand{\url}[1]{\hrefsym{#1}{\nolinkurl{#1}}}



% -- QR-Codes
\usepackage{qrcode} %-- hinter hyperref laden!

%--Römische Zahlen
\newcommand{\RM}[1]{\MakeUppercase{\romannumeral #1{}}}

%%--Abkürzungen etc.
\newcommand{\light}{\text{\Large $\lightning$}}

%-- Definitionen von weiteren Mathe-Befehlen
\DeclareMathOperator{\re}{Re} %Realteil
\DeclareMathOperator{\im}{Im} %Imaginaerteil
\DeclareMathOperator{\id}{id} %identische Abbildung
\DeclareMathOperator{\Sp}{Sp} %Spur
\DeclareMathOperator{\sgn}{sgn} %Signum
\DeclareMathOperator{\alt}{Alt} %Alternierende n-linearFormen
\DeclareMathOperator{\End}{End} %Endomorphismen
\DeclareMathOperator{\Vol}{Vol} %Volumen
\DeclareMathOperator{\dom}{dom} %Domain
\DeclareMathOperator{\Hom}{Hom} %Homomorphismen
\DeclareMathOperator{\Iso}{Iso} %Isomorphismen
\DeclareMathOperator{\END}{End} %Isomorphismen
\DeclareMathOperator{\bild}{Bild} %Bild
\DeclareMathOperator{\Kern}{Kern}
\DeclareMathOperator{\ev}{ev} %Auswertungsabbildung
\DeclareMathOperator{\rg}{rg} %Rang
\DeclareMathOperator{\Rg}{Rg} %Rang
\DeclareMathOperator{\diam}{diam} %Durchmesser
\DeclareMathOperator{\dist}{dist} %Distanz
\DeclareMathOperator{\grad}{grad} %Gradient
\DeclareMathOperator{\dive}{div} %Gradient
\DeclareMathOperator{\rot}{rot} %Rotation
\DeclareMathOperator{\hess}{Hess} %Hesse-Matrix
\DeclareMathOperator{\koker}{Koker} %Kokern
\DeclareMathOperator{\coker}{coker} %Kokern
\DeclareMathOperator{\aut}{Aut} %Automorphismen
\DeclareMathOperator{\ord}{ord} %Ordnung
\DeclareMathOperator{\ggT}{ggT} %ggT
\DeclareMathOperator{\kgV}{kgV} %kgV
\DeclareMathOperator{\Gr}{Gr} %Gerade
\DeclareMathOperator{\Kr}{Kr} %Kreis
\DeclareMathOperator{\Char}{char} %Charakteristik
\DeclareMathOperator{\Aut}{Aut} %Automorphismen
% \DeclareMathOperator{\D}{D} %formale Ableitung
\newcommand{\D}{\ensuremath{\mathrm{D}\mkern-1.5mu}}
\newcommand{\mathd}{\ensuremath{\mathrm{d}\mkern-1.0mu}}
\newcommand{\Tmap}{\ensuremath{\mathrm{T}\mkern-0.85mu}}
\newcommand{\opL}{\ensuremath{\mathrm{L}\mkern-0.6mu}}
\DeclareMathOperator{\cov}{cov} %Kovarianz
\DeclareMathOperator{\Gal}{Gal} %Galois
\DeclareMathOperator{\supp}{supp} %Träger
\DeclareMathOperator{\Sym}{Sym} %Symmetrische Gruppe
\DeclareMathOperator{\Zyl}{Zyl} %Zylinder
\DeclareMathOperator{\Mod}{Mod} %Moduln
\DeclareMathOperator{\EPK}{EPK} %Einpunktkompaktifizierung
\DeclareMathOperator{\conj}{conj} %Konjugation
\DeclareMathOperator{\ann}{ann} %Annulator
\DeclareMathOperator{\tor}{tor} %Torsionsmodul
\DeclareMathOperator{\Int}{Int} %Interior/Inneres
\DeclareMathOperator{\Br}{Br} %Brauerergruppe
\DeclareMathOperator{\GL}{GL} %allgemeine lineare Gruppe
\DeclareMathOperator{\SL}{SL} %spezielle lineare Gruppe
\DeclareMathOperator{\SU}{SU} %spezielle unitäre Gruppe
\DeclareMathOperator{\Tr}{Tr} %Spur/Trace
\DeclareMathOperator{\conv}{conv} %Konvexe Teilmenge
\DeclareMathOperator{\Span}{span} %span
\DeclareMathOperator{\Diff}{Diff} %Menge der Diffeomorphismen
\DeclareMathOperator{\Isom}{Isom} %Menge der Isometrien
\DeclareMathOperator{\spec}{spec} %Spektrum
\DeclareMathOperator{\res}{res} %Resolventenmenge
\DeclareMathOperator{\resolv}{R} %Resolventenmenge
\DeclareMathOperator{\idx}{ind} %Fredholm Index
\DeclareMathOperator{\Fred}{Fred} %menge der Fredholm-Operatoren

\newcommand{\tr}{\mathrm{tr}}
\newcommand{\w}{\mathrm{w}}
\newcommand{\sa}{\mathrm{s.a.}}
\newcommand{\so}{\mathrm{s.o.}}
\newcommand{\weakT}[1]{\ensuremath{\mathcal{T}_{#1}^{\mkern+1.0mu\text{\raisebox{0.4ex}{$\mathrm{w}$}}}}}
\newcommand{\weakTstar}[1]{\ensuremath{\mathcal{T}_{#1}^{\mkern+1.0mu\text{\raisebox{0.4ex}{$\mathrm{w}$}}^*}}}
\newcommand{\finSub}{\subset\mkern-0.7mu \subset}
\newcommand{\rand}[1]{\ensuremath{\partial^{\scriptscriptstyle #1}}}


% -- Kategorien
\DeclareMathOperator{\Ob}{Ob} %Objekte
\DeclareMathOperator{\obj}{obj} %Objekte
\DeclareMathOperator{\Mor}{Mor} %Morphismen
\newcommand{\stiso}{\mathrm{st. iso}} % stabiler Isomorphismus
\newcommand{\st}{\mathrm{st}}

\newcommand{\TOP}{\textsc{Top}}
\newcommand{\HTOP}{\textsc{HTop}}
\newcommand{\VR}{\textsc{VR}}
\newcommand{\MOD}{\textsc{Mod}}
\newcommand{\MONOIDE}{\textsc{Monoide}}
\newcommand{\MENGEN}{\textsc{Mengen}}
\newcommand{\MAN}{\textsc{Man}}
\newcommand{\GRUPPEN}{\textsc{Gruppen}}
\newcommand{\ABELGRUPPEN}{\textsc{Abel.Gruppen}}
\newcommand{\KAT}{\textsc{Kat}}
\newcommand{\FUN}{\textsc{Fun}}
\newcommand{\SIMP}{\textsc{Simp}}
\newcommand{\VEKT}{\textsc{Vekt}}


\newcommand{\op}{\mathrm{op}}
\newcommand{\ab}{\mathrm{ab}}
\newcommand{\pr}{\mathrm{pr}}
\newcommand{\pt}{\mathrm{pt}}
\newcommand{\const}{\mathrm{const}}
\newcommand{\CW}{\ensuremath{\mathrm{CW}}}
\DeclareMathOperator{\cell}{cell}
\DeclarePairedDelimiter{\homologieklasse}{\llbracket}{\rrbracket}

\newcommand{\Ker}{\ensuremath{\Kern \,}}
\newcommand{\zirkel}{\rotatebox[origin = c]{-90}{$\varangle$}}
\newcommand{\wirkung}{\!\curvearrowright\!}

\newcommand{\oE}{\OE}


%--Beweisende
\newcommand{\bewende}{\ifmmode \tag*{$\square$} \else \hfill $\square$ \fi}


% -- theorem packages
\usepackage{amsthm}
\usepackage{thmtools,thm-restate}
\renewcommand{\listtheoremname}{Übersicht aller Aussagen}

% -- Theoreme als PDF-Lesezeichen
\usepackage{bookmark}
\bookmarksetup{open,numbered}
\makeatletter
\newcommand*{\theorembookmark}{%
  \bookmark[
    dest=\@currentHref,
    rellevel=1,
    keeplevel,
  ]{%
    \thmt@thmname\space\csname the\thmt@envname\endcsname
    \ifx\thmt@shortoptarg\@empty
    \else
      \space(\thmt@shortoptarg)%
    \fi
  }%
}   
\makeatother

% -- Definition der einzelnen Umgebungen
\declaretheoremstyle[%
	headfont=\sffamily\bfseries,
	notefont=\normalfont\sffamily\scshape,
	bodyfont=\normalfont,
	headformat=\NUMBER\ \NAME\NOTE,
	headpunct=.,
	postheadspace=1em,
	spaceabove=15pt,spacebelow=10pt,
	postheadhook=\theorembookmark]%
{mainstyle}
\declaretheoremstyle[%
	headfont=\sffamily\bfseries,
	notefont=\normalfont\sffamily\scshape,
	bodyfont=\normalfont,
	headformat=\NUMBER\NAME\NOTE,
	headpunct=.,
	postheadspace=1em,
	spaceabove=15pt,spacebelow=10pt,
	postheadhook=\theorembookmark]%
{mainstyle_unnumbered}
\declaretheoremstyle[%
	headfont=\sffamily\bfseries,
	notefont=\normalfont\sffamily\scshape,
	bodyfont=\normalfont,
	headformat=swapnumber,
	headpunct=.,
	postheadspace=1em,
	spaceabove=15pt,spacebelow=10pt,
	postheadhook=\theorembookmark,
	qed=\qedsymbol]%
{mainstyleB}
\declaretheorem[name=Definition,parent=section,style=mainstyle]{definition}
\declaretheorem[name=Definition,numbered=no,style=mainstyle_unnumbered]{definition*}
\declaretheorem[name=Theorem,sharenumber=definition,style=mainstyle]{theorem}
\declaretheorem[name=Theorem,numbered=no,style=mainstyle_unnumbered]{theorem*}
\declaretheorem[name=Proposition,sharenumber=definition,style=mainstyle]{proposition}
\declaretheorem[name=Lemma,sharenumber=definition,style=mainstyle]{lemma}
\declaretheorem[name=Satz,sharenumber=definition,style=mainstyle]{satz}
\declaretheorem[name=Korollar,sharenumber=definition,style=mainstyle]{korollar}
\declaretheorem[name=Korollar,sharenumber=definition,style=mainstyleB]{korollarB}


% -- Beweise
\declaretheoremstyle[headfont=\sffamily\itshape\scshape,bodyfont=\normalfont,headpunct=:,postheadspace=1em,spacebelow=12pt,spaceabove=2pt,qed=\qedsymbol]{beweise}
\declaretheoremstyle[headfont=\sffamily\bfseries,bodyfont=\normalfont,headpunct=:,postheadspace=1em,spacebelow=10pt,spaceabove=10pt]{bemerkungen}
\declaretheorem[name=Beweis,numbered=no,style=beweise]{beweis}
\declaretheorem[name=Bemerkung,numbered=no,style=bemerkungen]{bemerkung}
\declaretheorem[name=Beispiel,numbered=no,style=bemerkungen]{beispiel}
\declaretheorem[name=Übung,numbered=no,style=bemerkungen]{uebung}
\declaretheorem[name=Erinnerung,numbered=no,style=bemerkungen]{erinnerung}

%--Volumen
\newcommand{\vol}[1]{\ensuremath{\Vol_{#1}}}

%--Skalarprodukt
\DeclarePairedDelimiterX\sprod[2]{\langle}{\rangle}{#1\,\delimsize\vert\,#2}
\DeclarePairedDelimiterX\skal[2]{\langle}{\rangle}{#1\,,\,#2}

%--Betrag, Gaußklammer
\DeclarePairedDelimiter{\abs}{\lvert}{\rvert}
\DeclarePairedDelimiter{\floor}{\lfloor}{\rfloor}
\DeclarePairedDelimiter{\ceil}{\lceil}{\rceil}

%--Norm
\DeclarePairedDelimiter\doppelstrich{\Vert}{\Vert}
\newcommand{\norm}[2][\relax]{
\ifx#1\relax \ensuremath{\doppelstrich*{#2}}
\else \ensuremath{\doppelstrich*{#2}_{#1}}
\fi}


%--Umklammern
\DeclarePairedDelimiter\enbrace{(}{)}
\DeclarePairedDelimiter\benbrace{[}{]}
\DeclarePairedDelimiter\lenbrace{<}{>}
\newcommand{\ssbrace}[1]{{\scriptscriptstyle\enbrace{#1}}}

%--Mengen
\DeclarePairedDelimiterX\mengenA[1]{\lbrace}{\rbrace}{#1}
\DeclarePairedDelimiterX\mengenB[2]{\lbrace}{\rbrace}{#1\, \delimsize\vert \, #2}

\makeatletter
\newcommand{\set}[2][\relax]{
\ifx#1\relax \ensuremath{
\mengenA*{#2}}
\else \ensuremath{%
  \mengenB*{#1}{#2}}
\fi}
\makeatother

\DeclareRobustCommand{\minwidthbox}[2]{%
  \ifmmode
    \expandafter\mathmakebox
  \else
    \expandafter\makebox
  \fi
  [\ifdim#2<\width\width\else#2\fi]{#1}%
}


%--offen und abgeschlossen
\newcommand{\off}{\! \stackrel[\text{offen}]{}{\subset} \!}
\newcommand{\abg}{\! \stackrel[\text{abg.}]{}{\subset} \!}

%--Differential
\newcommand{\diff}[2]{\ensuremath{\frac{{\partial #1}}{{\partial #2}} }}
\newcommand{\diffd}[2]{\ensuremath{\frac{\mathd #1}{\mathd #2} }}

%--Inhaltsverzeichnis
\usepackage[tocindentauto]{tocstyle}
\usetocstyle{KOMAlike}	
\shorthandon{"}
		
%--Punkte (als letztes laden)
\usepackage{ellipsis}

\newcommand{\fach}{$K$-Theorie und die Hopf-Invariante}
\newcommand{\semester}{SoSe 2015}
\newcommand{\homepage}{https://wwwmath.uni-muenster.de/reine/u/topos/lehre/SS2015/KTheorie-Hopf/Hopf.html}

\newcommand{\prof}{Dr.\ Ulrich Penning}
\publishers{\scalebox{14}{$K$}}
\providecommand{\verfasser}{Jannes Bantje}

%--Konfiguration von scrheadings
\setheadsepline{1pt}[\color{light_gray}]
\pagestyle{scrheadings}
%\fancyhf[H, F]{}
\clearscrheadfoot

\providecommand{\shortFach}{\fach}
\lehead{ \includegraphics[height=0.5 cm,keepaspectratio]{../!config/Bilder/Logo_WWU_Muenster_light_gray.pdf}}
\rehead{\rule{0cm}{0.5cm}\footnotesize \sffamily \color{light_gray} \verfasser -- Mitschrift \shortFach}
\lohead{\rule{0cm}{0.5cm} \footnotesize \sffamily \color{light_gray} Stand: \today \; \thistime[:]}
\rohead{\includegraphics[height=0.5 cm,keepaspectratio]{../!config/Bilder/fb10logo_gray.pdf}}


\ofoot[{ \color{dark_gray} \LARGE \sffamily \thepage}]{{ \color{dark_gray} \LARGE \sffamily \thepage}} %hier wir auch der plain Stil bearbeitet!
\automark{section}
\ifoot{ \color{dark_gray} \small \leftmark}

%--Metadaten
\author{\verfasser}
\titlehead{\includegraphics[height=1.5cm, keepaspectratio]{../!config/Bilder/Logo_WWU_Muenster.pdf}%
\hfill \includegraphics[height=1.3cm, keepaspectratio]{../!config/Bilder/fb10logo.pdf}}
\title{Skript \fach}
\subtitle{Mitschrift der Vorlesung  \enquote{\fach} von \prof}

\publishers{
Aktuelle Version verfügbar bei: \bigskip\\
\normalsize
\begin{minipage}[t]{0.45\textwidth}
	\centering{
	\qrcode[height=3.5cm, version=4]{https://github.com/JaMeZ-B/latex-wwu}
	\medskip\\
	\includegraphics[height=0.4cm, keepaspectratio]{../!config/Bilder/github_octo.pdf} 
	\includegraphics[height=0.4cm, keepaspectratio]{../!config/Bilder/GitHub_Logo.pdf} (inklusive Sourcecode)\\
	{\small\url{https://github.com/JaMeZ-B/latex-wwu}}}
\end{minipage}
\hfill
\begin{minipage}[t]{0.45\textwidth}
	\centering{
	\qrcode[height=3.5cm, version=4]{btsync://B6WH2DISQ5QVYIRYIEZSF4ZR2IDVKPN3I?n=latex_share}
	\medskip\\
	\includegraphics[height=0.4cm, keepaspectratio]{../!config/Bilder/bt_sync_logo.pdf} 
	{\large \textbf{Bittorrent} Sync} \\
	\texttt{B6WH2DISQ5QVYIRYIEZSF4ZR2IDVKPN3I}
	}
\end{minipage}
}


\begin{document}
\pagenumbering{Roman}
\maketitle
\begin{abstract}
\section*{Vorwort --- Mitarbeit am Skript}
Dieses Dokument ist eine Mitschrift aus der Vorlesung \enquote{\fach, \semester}, gelesen von \prof. Der Inhalt entspricht weitestgehend dem Tafelanschrieb. Für die
Korrektheit des Inhalts übernehme ich keinerlei Garantie! Für Bemerkungen und Korrekturen -- und seien es nur Rechtschreibfehler -- bin ich sehr dankbar. 
Korrekturen lassen sich prinzipiell auf drei Wegen einreichen: 
\begin{itemize}
	\item Persönliches Ansprechen in der Uni, Mails an \hrefsymmail{mailto:\mail}{\mail} (gerne auch mit annotieren PDFs) 
	\item \emph{Direktes} Mitarbeiten am Skript: Den Quellcode poste ich auf GitHub (siehe oben), also stehen vielfältige Möglichkeiten der Zusammenarbeit zur Verfügung:
	Zum Beispiel durch Kommentare am Code über die Website und die Kombination Fork + Pull Request. Wer sich verdient macht oder ein Skript zu einer Vorlesung, die 
	ich nicht besuche, beisteuern will, dem gewähre ich gerne auch Schreibzugriff.
	
	Beachten sollte man dabei, dass dazu ein Account bei \url{github.com} notwendig ist, der allerdings ohne Angabe von persönlichen Daten angelegt werden kann. 
	Wer bei GitHub (bzw. dem zugrunde liegenden Open-Source-Programm \enquote{\texttt{git}}) -- verständlicherweise -- Hilfe beim Einstieg braucht, dem helfe ich gerne 
	weiter. Es gibt aber auch zahlreiche empfehlenswerte Tutorials im Internet.\footnote{zB. \url{https://try.github.io/levels/1/challenges/1}, ist auf Englisch, aber dafür 
	interaktives LearningByDoing}
	\item \emph{Indirektes} Mitarbeiten: \TeX-Dateien per Mail verschicken. 
	
	Dies ist nur dann sinnvoll, wenn man einen ganzen Abschnitt ändern möchte (zB. einen alternativen Beweis geben), da ich die Änderungen dann per Hand einbauen muss! Ich freue mich aber auch über solche Beiträge!
\end{itemize}
\section*{Vorlesungshomepage}
{\centering 
\begin{minipage}[c][][c]{\textwidth}
	\centering \qrcode[height=3cm]{\homepage} \medskip\\
	\footnotesize \url{\homepage}
\end{minipage}
\par}

\end{abstract}

\tableofcontents
\cleardoubleoddemptypage

\pagenumbering{arabic}
\setcounter{page}{1}

\setcounter{section}{-1}
\section{Einführung} % (fold)
\label{sec:0}

\subsection[Definition: Reelle Divisionsalgebra]{Definition} % (fold)
\label{sub:01}
Eine $\mathds{R}$-Algebra $\mathcal{A}$ mit $1$, die nicht notwendigerweise assoziativ ist, heißt \Index{reelle Divisionsalgebra}, falls jedes Element $a\not= 0$ invertierbar 
ist.
\minisec{Beispiele}
Die reellen Zahlen $\mathds{R}$, die komplexen Zahlen $\mathds{C}$, die Quaternionen\footnote{siehe auch \url{https://de.wikipedia.org/wiki/Quaternion}} $\mathds{H}$ und die sogenannten Cayley-Zahlen\footnote{auch Oktonionen oder reelle Okataven; siehe  \url{http://de.wikipedia.org/wiki/Oktave_(Mathematik)}} $\mathds{O}$. Dabei ist
$\mathds{O} \cong \mathds{H} \oplus \mathds{H}$, wobei die Multiplikation wie folgt definiert ist:
\[
	(a,b) \cdot (d,c) = (ac - d^* b, da + bc^*) 
\]
Dabei ist $d= x + iy + jz + kw$ und $d^* = x -iy -jz - kw$.
% subsection 11 (end)

\subsection{Theorem 1.1 (Adams)} % (fold)
\label{sub:02}
Folgende Aussagen sind äquivalent:
\begin{enumerate}[(i)]
	\item $\mathds{R}^n$ besitzt die Struktur einer reellen Divisionsalgebra.
	\item Entweder ist $n=1$ oder $n \ge 2$ ist gerade und es gibt eine stetige Abbildung $f \colon S^{2n-1} \to S^n$ mit der sogenannten \Index{Hopf-Invariante} $1$.
	\item Es gilt $n \in \set{1,2,4,8}$.
\end{enumerate}
% subsection 11 (end)

\subsection{Geplanter Verlauf der Vorlesung} % (fold)
\label{sub:inhalt}
Wir werden diesen Satz mit Hilfe der sogenannten $K$-Theorie beweisen. Dazu müssen wir aber einiges an Vorarbeit leisten:
\begin{enumerate}[1.]
	\item Verallgemeinerte Kohomologietheorien
	\item Vektorbündel
	\item $K$-Theorie
	\begin{enumerate}[label*=\arabic*.]
		\item Produkte in $K$-Theorie
		\item Bott-Periodizität
		\item Thom-Isomorphismus
	\end{enumerate}
	\item Hopf-Invariante und der Beweis des Satzes
\end{enumerate}
% subsection inhalt (end)
% section 0 (end)
\newpage

\section{Kohomologietheorien} % (fold)
\label{sec:1}

\subsection[Definition: Kontravarianter Funktor]{Definition} % (fold)
\label{sub:11}
Seien $\mathcal{C}$ und $\mathcal{D}$ Kategorien. Ein \Index{kontravarianter Funktor} $F \colon \mathcal{C} \to \mathcal{D}$ 
\begin{itemize}
	\item ordnet jedem Objekt $c \in \mathrm{Obj}(\mathcal{C})$ ein Objekt $F(c) \in \mathrm{Obj}(\mathcal{D})$ zu
	\item ordnet jedem Morphismus $f \colon c \to c'$ in $\mathrm{Mor}_{\mathcal{C}}(c,c')$ einen Morphismus $F(f) \colon F(c') \to F(c)$ zu 
\end{itemize}
sodass folgende Eigenschaften gelten:
\begin{itemize}
	\item Für $f \colon c \to c'$ und $g \colon c' \to c''$ gilt $F(g \circ f) = F(f) \circ F(g)$
	\item $F(\id_c) = \id_{F(c)}$
\end{itemize}
\minisec{Beispiel (Dualraumfunktor)}
Sei $K$ ein Körper. Sei $\VEKT_K$ die Kategorie der endlich dimensionalen $K$-Vektorräume und linearen Abbildungen. Sei 
${ }^* \colon \VEKT_K \to \VEKT_K$ gegeben durch $V \mapsto \Hom_K(V,K)=V^*$. Für $f \colon V \to W$ sei $f^* \colon W^* \to V^*$, 
$\varphi \mapsto \varphi \circ f$. Dies ist ein kontravarianter Funktor. Insbesondere ist $(f \circ g)^* = g^* \circ f^*$.
% subsection 11 (end)

\subsection[Definition: Verallgemeinerte Kohomologietheorie]{Definition} % (fold)
\label{sub:12}
Sei $\TOP^2$ die Kategorie der Raumpaare $(X,A)$. Sei $V \colon \TOP^2 \to \TOP^2$ der Funktor $(X,A) \mapsto(A,\emptyset)$. Sei $R$ ein kommutativer Ring mit $1$. Eine 
\Index{verallgemeinerte Kohomologietheorie} $h^* = (h^n, \partial^n)_{n \in \mathds{Z}}$ mit Werten in $R$-Moduln ist ein kontravarianter Funktor 
$h^* \colon \TOP^2 \to \textsc{Gr}\text{-}R\text{-}\MOD$ zusammen mit einer natürlichen Transformation 
\[
	\partial^* \colon h^* \circ V \to h^{*+1},
\]
so dass die folgenden Eigenschaften gelten:
\begin{enumerate}[a)]
	\item \Index{Homotopieinvarianz}: Seien $f,g \colon (X,A) \to (Y,B)$ homotope Abbildungen von Raumpaaren. Dann gilt für alle $n \in \mathds{Z}$: 
	\[
		h^n(f)= h^n(g)
	\]
	\item \Index{Lange exakte Paarsequenz}: Sei $(X,A)$ ein Raumpaar. Seien $(A,\emptyset) \xhookrightarrow{i}(X,\emptyset)$ und $j \colon (X,\emptyset) \to (X,A)$ die 
	kanonischen Inklusionen. Dann ist die Folge 
	\[
		\begin{tikzcd}
			\cdots \rar["\partial^{n-1}"] & h^n(X,A;R) \rar["h^{n}(j)"] & h^n (X,\emptyset;R) \rar["h^n(i)"] & h^n(A, \emptyset;R) \rar["\partial^n"] & h^{n+1}(X,A;R) \rar & 
			\cdots
		\end{tikzcd}
	\]
	\item \Index{Ausschneidung}: Sei $(X,A)$ ein Raumpaar, $U \subset A$, sodass $\overline{U} \subseteq \mathring{A}$. Dann ist die von der Inklusion 
	$\iota \colon (X \setminus U, A \setminus U) \hookrightarrow (X,A)$ induzierte Abbildung $h^n(\iota) \colon h^n(X,A;R) \to h^n(X \setminus U, A \setminus U;R)$ ein Isomorphismus.
\end{enumerate}
\minisec{Bemerkungen}
\begin{enumerate}[(i)]
	\item Wir schreiben kurz $h^n(X,A)$ für $h^n(X,A;\mathds{Z})$ und $f^*$ für $h^n(f)$, falls sich der Grad aus dem Kontext ergibt. 
	\item $h^n(\set{\mathrm{pt}}) =: h^n$ heißen \Index{Koeffizienten der Kohomologietheorie}.
	\item Einige Kohomologietheorien besitzen ein \Index{externes Produkt}, das heißt eine natürliche Transformation\todo{Diagramm hinzufügen}
	\[
		h^n(X,A;R) \times h^m(Y,B;R) \longrightarrow h^{n+m} \enbrace[\big]{(X,A) \times (Y,B);R} 
	\]
	die bilinear und assoziativ ist. 
	Dabei ist $(X,A) \times (Y,B) = \enbrace[\big]{X \times Y, (A \times Y) \cup (X \times B)}$.
	
	Falls ein externes Produkt für $h^*$ existiert, dann ist $h^*(X,A;R) = \bigoplus_{n \in \mathds{Z}} h^n(X,A;R)$ ein \Index{graduierter Ring} mit der Multiplikation 
	\[
		\begin{tikzcd}[row sep=large, column sep=small]
			h^n(X,A;R) \times h^m(X,A;R) \arrow[dashed]{rr} \arrow["\times"']{dr} & & h^{n+m}(X,A;R) \\
			& h^{n+m} \enbrace[\big]{(X,A) \times (X,A);R} \arrow["h^{n+m}(\Delta)"']{ur} &
		\end{tikzcd}
	\]
\end{enumerate}
% subsection 12 (end)

\subsection[Definition: Ko-Raumpaar und Kofaserung]{Definition} % (fold)
\label{sub:13}
Ein Raumpaar $(X,A)$ heißt \Index{Ko-Raumpaar}, falls
\begin{itemize}
	\item $X$ ein kompakter Hausdorffraum ist
	\item $A \subset X$ abgeschlossen ist
	\item die Inklusion $A \hookrightarrow X$ eine \Index{Kofaserung} ist, das heißt für jede stetige Abbildung\marginnote{Homotopieerweiterungseigenschaft} 
	$f \colon A \times I \cup X \times \set{0} \to Y$ existiert 
	$H \colon X \times I \to Y$, sodass das folgende Diagramm kommutiert:
	\[
		\begin{tikzcd}[sep=large]
			A \times I \cup X \times \set{0} \rar["f"] \dar[hook]& Y \\
			X \times I \urar["H"'] & 
		\end{tikzcd}
	\] 
\end{itemize}
Mit $\TOP^2_{\mathrm{Kof}} \subseteq \TOP^2$ bezeichnen wir die Kategorie der Ko-Raumpaare und stetigen Abbildungen.
\minisec{Beispiel}
\begin{itemize}
	\item Sei $X$ ein kompakter \CW-Komplex, $A \subset X$ ein Unterkomplex. Dann ist $\iota \colon A \hookrightarrow X$ eine Kofaserung.
	\item $X$ eine kompakte Mannigfaltigkeit, $A$ eine abgeschlossene Untermannigfaltigkeit.
\end{itemize}
\minisec{Bemerkung}
Falls $(X, \set{x_0})$ ein Objekt in $\TOP^2_{\mathrm{Kof}}$ ist, dann heißt $X$ \Index{wohlpunktiert}.
\minisec{Beispiel}
$X= [-1,1] \times [-1,1]$, $x_0=(0,0)$, $Y=X$. Sei $\gamma \colon I \to X$ ein Pfad mit $\gamma(0)= (0,0)=x_0$. Wir erhalten eine Abbildung 
$f \colon \set{x_0} \times I \cup  X \times \set{0} \to X $
\missingfigure{Skizze zu der Abbildung}
Kofaserung heißt in diesem Fall, dass der Raum $X$ entlang des Pfades \enquote{mitgezogen} werden kann.
% subsection 13 (end)
% section 1 (end)
\newpage

\section{Vektorbündel} % (fold)
\label{sec:2}
Sei für den Verlauf dieses Kapitels $K \in \set{\mathds{R},\mathds{C}}$. 

\subsection[Definition: $K$-Vektorbündel]{Definition} % (fold)
\label{sub:21}
Sei $X$ ein topologischer Raum. Eine stetige surjektive Abbildung $p \colon E \to X$ heißt $K$-\Index{Vektorbündel}, falls gilt: 
\begin{enumerate}[a)]
	\item $p^{-1}(\set{x}) =: E_x$ ist ein $K$-Vektorraum für alle $x \in X$. $E_x$ heißt die \Index{Faser} von $E$ über $x$. 
	\item Für alle $x \in X$ gibt es eine Umgebung $U$, ein $n \in \mathds{N}$ und einen Homöomorphismus $\varphi_U \colon p^{-1}(U) \to U \times K^n$, sodass das folgende 
	Diagramm kommutiert
	\[
		\begin{tikzcd}
			p^{-1}(U) \arrow["\varphi_U"]{rr} \drar["p"'] & & U \times K^n  \dlar["\mathrm{pr}_U"] \\
			& U 
		\end{tikzcd}
	\]
	und die Einschränkung $\varphi_U\big|_{E_x} \colon E_x \to K^n$ für alle $x \in U$ ein Isomorphismus von $K$-Vektorräumen ist. Diese Eigenschaft wird oft als 
	\Index{lokale Trivialität} bezeichnet.
\end{enumerate}
\minisec{Beispiel}
\begin{itemize}
	\item Sei $X$ ein beliebiger topologischer Raum. Dann ist $X \times K^n \xrightarrow{\mathrm{pr}_X} X$ das \Index{triviale Vektorbündel} über $X$.
	\item Betrachte $S^n \subset \mathds{R}^{n+1}$. Sei $\Tmap S^n = \set[(x,v) \in S^n \times \mathds{R}^{n+1}]{\skal*{x}{v}=0}$. Definiere $p \colon \Tmap S^n \to S^n$ 
	durch $(x,v) \mapsto x$. Damit erhalten wir ein $\mathds{R}$-Vektorbündel über $S^n$ (Tangentialbündel). 
	
	\emph{Übung: Zeige lokale Trivialität. Siehe Anhang \ref{sub:lokal_trivial_Sn}}
\end{itemize}
\minisec{Bemerkung}
Die Abbildung $x \mapsto \dim_K (E_x)$ ist stetig, also lokal konstant, das heißt über einem zusammenhängenden Raum hat ein Vektorbündel konstante Faserdimension. 
% subsection 21 (end)

\subsection[Definition: Vektorraumbündelmorphismus]{Definition} % (fold)
\label{sub:22}
Seien $p \colon E \to X$ und $p' \colon E' \to X'$ zwei $K$-Vektorbündel. Ein \bet{Morphismus von Vektorbündeln}\index{Morphismus!von Vektorbündeln} oder 
\Index{Vektorraumbündelmorphismus} von $E$ nach $E'$ ist ein Paar $(f, \overline{f})$ von stetigen Abbildungen $\overline{f} \colon E \to E'$ und $f \colon X \to X'$, sodass 
das folgende Diagramm kommutiert
\[
	\begin{tikzcd}[sep=large]
		E \rar["\overline{f}"] \dar["p"'] & E' \dar["p'"] \\
		X \rar["f"] & X'
	\end{tikzcd}
\]
und sodass die Einschränkung $\overline{f}\big|_{E_x} \colon E_x \to E'_{f(x)}$ für jedes $x \in X$ ein $K$-Vektorraumhomomorphismus ist. Wir bezeichnen die Kategorie der 
endlich dimensionalen $K$-Vektorraumbündel über einem topologischen Raum $X$ zusammen mit den Morphismen $(\id_X, \overline{f})$ mit $\VEKT_K(X)$.
% subsection 22 (end)

\section*{Operationen mit Vektorbündeln}
\subsection[Definition: Pullback]{Definition} % (fold)
\label{sub:23}
Seien $X$ und $Y$ topologische Räume. Sei $f \colon Y \to X$ eine stetige Abbildung und sei $p \colon E \to X$ ein Vektorbündel. Dann heißt
\[
	f^* E := \set[(y,v) \in Y \times E]{f(y)=p(v)\rule{0cm}{0.9em}} 
\]
zusammen mit der Abbildung $q \colon f^* E \to Y$, $(y,v) \mapsto y$ der \Index{Pullback}\marginnote{Der Pullback lässt sich auch allgemein auf kategorieller Ebene 
definieren} von $p \colon E \to X$ entlang von $f$. Mit $\overline{f}(y,v)=v$ kommutiert das 
folgende Diagramm
\[
	\begin{tikzcd}
		f^* E \rar["\overline{f}"] \dar["q"'] & E \dar["p"] \\
		Y \rar["f"] & X
	\end{tikzcd}
\]
% subsection 23 (end)

\subsection[Lemma: Universelle Eigenschaft des Pullbacks]{Lemma} % (fold)
\label{sub:24}
\begin{enumerate}[a)]
	\item $q \colon f^* E \to Y$ ist ein Vektorbündel und $(\overline{f},f)$ ist ein Morphismus von Vektorbündeln.
	\item $q \colon f^* E \to Y$ ist durch die folgende universelle Eigenschaft eindeutig (bis auf Isomorphie) charakterisiert:
	
	Sei $r \colon F \to Z$ ein Vektorbündel, $(\overline{g},g)$ ein Morphismus zwischen $r \colon F \to Z$ und $p \colon E \to X$. Sei weiter $g' \colon Z \to Y$ eine 
	stetige Abbildung mit $f \circ g' = g$. Dann existiert genau ein $G \colon F \to f^* E$, so dass $(G,g')$ ein Morphismus von Vektorbündeln ist und weiter 
	$(\overline{f},f) \circ (G,g')= (\overline{g},g)$ gilt.
	\[
		\begin{tikzcd}[sep=large]
			F \drar["\exists!\, G", dashed] \dar["r"'] \arrow[bend left,"\overline{g}"]{drr} & & \\
			Z \drar["g'"'] \arrow["g" near end]{drr} & f^* E \rar["\overline{f}"] \dar["q" pos=0.3,crossing over] & E \dar["p"]\\
			& Y \rar["f"'] & X
		\end{tikzcd}
	\]
\end{enumerate}
\minisec{Beweis}
Sei $y\in Y$. Sei $U \subset X$ eine Umgebung von $f(y)$ in $X$, sodass eine lokale Trivialisierung $\varphi_U \colon p ^{-1}(U) \to U \times K^n$ existiert. $V:= f^{-1}(U)$ 
ist eine Umgebung von $y$ und es gilt
\[
	q ^{-1}(V) = \set[(y,v) \in V \times p ^{-1}(U)]{f(y)=p(v)} 
\]
Wir definieren nun 
\begin{align*}
	\psi_V \colon \,\,q ^{-1}(V) &\longrightarrow V \times K^n &  (y,v)&\longmapsto \enbrace[\big]{y, \pr_{K^n} \circ \varphi_U(v)} \\
	\kappa_V \colon V \times K^n &\longrightarrow q ^{-1}(V) &  (y,w) &\longmapsto \enbrace*{y, \varphi_U^{-1} \enbrace*{f(y),w} } 
\end{align*}
Dann gilt $\psi_V \circ \kappa_V=\id_{V \times K^n}$ und $\kappa_V \circ \psi_V = \id_{q^{-1}(V)}$. Außerdem kommutiert das Diagramm aus \ref{sub:21}. Also ist $\psi_V$ eine 
lokale Trivialisierung von $f^*E \to Y$ über $V$.
Der Beweis von b) ist eine (relativ einfache) Übungsaufgabe\todo{eventuell hinzufügen}.\bewende

\bigskip

\noindent Sei $\mathcal{C}$ eine (topologische) Kategorie. Hieraus lässt sich eine neue Kategorie $\mathcal{C}^\op$ definieren:
\[
	\obj(\mathcal{C}^\op) := \obj(\mathcal{C}) \quad , \quad \Mor_{\mathcal{C}^\op}(c,d) := \Mor_{\mathcal{C}}(d,c)
\]
mit der neuen Komposition
\[
	\begin{tikzcd}[row sep=0.1em]
		\bullet \mbox{ :}\hspace{-4em} & \Mor_{\mathcal{C}^\op}(c',c'') \times \Mor_{\mathcal{C}^\op}(c,c') \rar & \Mor_{\mathcal{C}^\op}(c,c') \\
		& (f,g) \rar[mapsto] & f \bullet g = g \circ f
	\end{tikzcd}
\]
(Kovariante) Funktoren $\mathcal{C}^\op \to \mathcal{D}$ entsprechen dann kontravarianten Funktoren $\mathcal{C}\to \mathcal{D}$. Sei $\VEKT_K = \VEKT_K(\pt)$ die Kategorie der 
endlichdimensionalen $K$-Vektorräume und linearen Abbildungen. Ein Funktor
\begin{equation}
	F \colon \Underbrace{\VEKT_K \times \ldots \times \VEKT_K}{r} \times \Underbrace{\VEKT_K^\op \times \ldots \times \VEKT_K^\op}{s} \longrightarrow \VEKT_K \label{eq:24}\tag{\#}
\end{equation}
heißt \bet{stetig}\index{stetiger Funktor}, falls die induzierte Abbildung
\begin{gather*}
	\Mor_{\VEKT_K}(V_1,V'_1) \times \ldots \times \Mor_{\VEKT_K}(V_r,V'_r) \times \Mor_{\VEKT_K^\op}(V_{r+1},V'_{r+1}) \times \ldots \times \Mor_{\VEKT_K^\op}(V_{r+s},V'_{r+s}) \\
	\longrightarrow \enspace\Mor_{\VEKT_K} \enbrace*{F(V_1, \ldots,V_r,V_{r+1}, \ldots , V_{r+s}), F (V'_1, \ldots , V'_r, V'_{r+1}, \ldots , V'_{r+s})} 
\end{gather*}
stetig ist (dies klappt nur für Kategorien, für die eine Topologie auf den Morphismen existiert!).
% subsection 24 (end)

\subsection[Lemma: Induzierter Funktor eines stetigen Funktors]{Lemma} % (fold)
\label{sub:25}
Sei $F$ ein stetiger Funktor wie in \eqref{eq:24} und sei $X$ ein topologischer Raum. Dann induziert $F$ einen Funktor
\[
	F_X \colon \Underbrace{\VEKT_K(X) \times \ldots \times \VEKT_K(X)}{r} \times \Underbrace{\VEKT_K(X)^\op \times \ldots \times \VEKT_K(X)^\op}{s} \longrightarrow \VEKT_K(X)
\]
der verträglich ist mit Pullbacks und $F_\pt=F$ erfüllt.
\minisec{Beweis}
Sei $m=r+s$. Definiere 
\(
	F_X \enbrace*{E^\ssbrace{1}, \ldots , E^\ssbrace{m}} = \coprod_{x \in X} F \enbrace*{E_x^\ssbrace{1}, \ldots , E_x^\ssbrace{m}}   
\)
als Menge. Sei $U_i$ eine offene Überdeckung von $X$, so dass lokale Trivialisierungen 
\[
	\varphi_i^\ssbrace{k} \colon (p^\ssbrace{k})^{-1}(U_i) \longrightarrow U_i \times K^{\ell_k}
\]
existieren. $F$ induziert Bijektionen
\[
	F_{U_i} \enbrace*{\varphi_i^\ssbrace{1}, \ldots , \varphi_i^\ssbrace{m}} \colon F_{U_i} \enbrace*{(p^\ssbrace{1})^{-1}(U_i), \ldots , (p^\ssbrace{m})^{-1}(U_i)}
	\longrightarrow U_i \times F \enbrace*{K^{\ell_1}, \ldots , K^{\ell_m}}   
\]
Jetzt gibt es eine eindeutige Topologie auf $F_X \enbrace*{E^\ssbrace{1}, \ldots , E^\ssbrace{m}}$, so dass die eben definierten Abbildungen 
$\psi_i = F_{U_i} \enbrace*{\varphi_i^\ssbrace{1}, \ldots , \varphi_i^\ssbrace{m}}$ Homöomorphismen werden, wobei $U_i \times F \enbrace*{K^{\ell_1}, \ldots , K^{\ell_m}}$ 
die Produkttopologie trägt: Eine Menge $V \subset F_X \enbrace*{E^\ssbrace{1}, \ldots , E^\ssbrace{m}}$ ist offen, falls sie die Vereinigung von Mengen der Form 
$\psi_i^{-1}(V_i)$ mit $V_i \subset U_i \times F \enbrace*{K^{\ell_1}, \ldots , K^{\ell_m}}$ offen ist.

Damit diese Definition konsistent ist, müssen wir folgendes prüfen: Sei $V \subset (U_i \cap U_j) \times F(K^{\ell_1}, \ldots , K^{\ell_m})$ offen. Wir haben zu zeigen, dass 
dann auch $\psi_j \circ \psi_i^{-1}(V)$ offen ist. Aber es gilt $\psi_j \circ \psi_i ^{-1}(x,v)= \enbrace[\big]{x, \psi_{ij}(x)(v)}$ für eine Abbildung
\[
	\psi_{ij} \colon U_i  \cap U_j \longrightarrow \End_K \enbrace*{F(K^{\ell_1}, \ldots , K^{\ell_m})} 
\]
Es gilt $\psi_{ij} = F \enbrace[\big]{\psi_{ij}^\ssbrace{1}, \ldots , \psi_{ij}^\ssbrace{m}}$ für \emph{stetige} Abbildungen 
$\psi_{ij}^\ssbrace{k} \colon U_i \cap U_j \to \End_K(K^{e_k})$. Da $F$ stetig ist, folgt, dass alle $\psi_{ij}$ stetig sind und folglich auch $\psi_j \circ \psi_i ^{-1}$. 
Somit ist $\psi_j \circ \psi_i ^{-1}(V)$ offen. Es ist damit klar, dass $F_X \enbrace*{E^\ssbrace{1}, \ldots , E^\ssbrace{m}} \to X$ ein Vektorbündel ist, denn wir haben lokale Trivialisierungen konstruiert. Ferner gilt per Definition für eine stetige Abbildung $f \colon Y \to X$
\[
	f^* F_X \enbrace*{E^\ssbrace{1}, \ldots , E^\ssbrace{m}} \cong F_Y \enbrace*{f^*E^\ssbrace{1}, \ldots , f^* E^\ssbrace{m}}  \bewende
\]
\minisec{Beispiele}
\begin{itemize}
	\item $\oplus \colon \VEKT_K \times \VEKT_K \to \VEKT_K$ ist stetig $\Rightarrow $ direkte Summe von Vektorbündeln
	\item $\otimes \colon \VEKT_K \times \VEKT_K \to \VEKT_K$ ist stetig $\Rightarrow $ Tensorprodukt von Vektorbündeln
	\item Dualisieren von Vektorräumen ${ }^* \colon \VEKT_K^\op \to \VEKT_K$ ist stetig $\Rightarrow $ duales Vektorbündel
	\item $\Hom \colon \VEKT_K \times \VEKT_K^\op \to \VEKT_K$, $(V,W) \mapsto \Hom_K(W,V)$ ist stetig $\Rightarrow$ Homomorphismenbündel $\Hom(E,F)$.
\end{itemize}
Seien $p \colon E \to X$ und $q \colon F \to X$ zwei $K$-Vektorbündel. Dann ist $E \times F \xrightarrow{(p,q)} X \times X$ auch ein Vektorbündel, wenn $E_X \times F_X$
die Vektorraumstruktur der äußeren direkten Summe trägt.
\minisec{Übungsaufgabe}
Sei $\Delta  \colon X \to X \times X$ die Diagonalabbildung. Zeige, dass $\Delta^*(E \times F) \cong E \oplus F$ als Vektorbündel über $X$.
% subsection 25 (end)
\section*{Homotopieinvarianz von Pullbacks}
\minisec{Erinnerung}
Ein topologischer Raum heißt \bet{normal}\index{normaler topologischer Raum}, falls sich zwei disjunkte abgeschlossene Mengen durch offene Mengen trennen lassen. 
\minisec{Bemerkung}
Jeder kompakte Hausdorffraum ist normal. \emph{Übungsaufgabe}\todo{In den Anhang?}

\subsection{Satz (Fortsetzungssatz von Tietze)} % (fold)
\label{sub:26}
Sei $X$ ein normaler Raum, $A \subset X$ abgeschlossen und $V$ ein endlich dimensionaler $K$-Vektorraum. Zu jeder stetigen Abbildung $f \colon A \to V$ existiert eine 
stetige Abbildung $F \colon X \to V$ mit $F|_A=f$.
\minisec{Beweis}

\minisec{Definition}
Sei $p \colon E \to X$ ein $K$-Vektorbündel über einem topologischen Raum $X$. Ein \Index{Schnitt} von $E$ ist eine stetige Abbildung $s \colon X \to E$ mit 
$p \circ s = \id_X$. Das heißt $s(x) \in E_x$. Jedes Vektorbündel hat einen Schnitt, nämlich den trivialen Schnitt $s(x)=0 \in E_x$.
% subsection 26 (end)

\subsection[Lemma über Fortsetzungen von Schnitten auf abgeschlossenen Teilmengen]{Lemma} % (fold)
\label{sub:27}
Sei $X$ ein kompakter Hausdorffraum, $A \subset X$ abgeschlossen. Seien $p \colon E \to X$ und $p' \colon E' \to X$ zwei $K$-Vektorbündel. Dann gilt:
\begin{enumerate}[a)]
	\item Jeder Schnitt $s \colon A \to E|_A$ lässt sich zu einem Schnitt $\overline{s} \colon X \to E$ fortsetzen.
	\item Sei $\iota \colon A \hookrightarrow X$ die Inklusion und $E|_A := \iota^* E$, sowie $E'|_A := \iota^* E'$. Sei weiter $f \colon E|_A \to E'|_A$ ein Morphismus von 
	Vektorbündeln. Dann existiert ein Morphismus $\hat{f} \colon E \to E'$, der $f$ fortsetzt. Ist $f$ ein Isomorphismus, dann existiert eine offene Menge $U \supset A$, 
	sodass $\hat{f}|_U \colon E|_U \to E'|_U$ ein Isomorphismus von Vektorbündeln ist.
\end{enumerate}
\minisec{Beweis}
\begin{enumerate}[a)]
	\item Wähle eine endliche Überdeckung $(U_i)_{i=1,\ldots ,N}$ von $X$, sodass lokale Trivialisierungen $\varphi_i \colon E|_{U_i} \to U_i \times K^n$ existieren. Nach 
	dem Fortsetzungssatz lässt sich 
	\[
		\pr_{K^n} \circ \varphi_i \circ s\big|_{U_i \cap A} \colon U_i \cap A \longrightarrow K^n 
	\]
	zu einer stetigen Abbildung auf $U_i$ fortsetzen, also auch $s|_{U_i \cap A}$ selbst. Sei $\overline{s}_i$ eine solche Fortsetzung. Sei $\psi_i \colon X \to [0,1]$ eine 
	Partition der Eins\footnote{siehe \url{https://de.wikipedia.org/wiki/Zerlegung_der_Eins}}, die der Überdeckung $(U_i)_{i=1,\ldots ,N}$ untergeordnet ist. Sei
	\[
		\overline{s} = \sum_{i=1}^{N} \psi_i \cdot \overline{s}_i \colon X \longrightarrow E 
	\]
	Dann gilt $\overline{s}|_A = \enbrace*{\sum_{i=1}^{N} \psi_i} \cdot s|_A= s|_A$.
	\item Die Morphismen $E \xrightarrow{f} E'$ entsprechen den Schnitten des Vektorbündels $\Hom(E,E') \to X$ wie folgt: Jedem $f$ lässt sich ein Schnitt
	\[
		s_f(x)(v) := f(v)
	\]
	zuordnen. Jedem Schnitt $s \colon X \to \Hom(E,E')$ entspricht ein Morphismus $f_s(v)=s \enbrace*{p(v)}(v) \in E'_{p(v)}$. Somit folgt der erste Teil aus a).
	
	Sei $\pi \colon \Hom(E,E') \to X$ die Bündelabbildung. Sei weiter
	\[
		\Iso(E,E') = \set[g \in \Hom(E,E')]{g \colon E_{\pi(g)} \to E'_{\pi(g)} \text{ ist ein Isomorphismus}} 
	\]
	Dann ist $\Iso(E,E') \to X$ ein lokal triviales Bündel, allerdings \emph{kein} Vektorbündel. Sind $E$ und $E'$ $n$-dimensional, ist die Faser von $\Iso(E,E')$ homöomorph
	zu $\GL_n(K)$. Da $\GL_n(K) \subset M_n(K)$ offen ist und $\Iso(E,E')$ lokal trivial ist, folgt, dass $\Iso(E,E') \subset \Hom(E,E')$ offen ist.\footnote{da lokal 
	homöomorph zu$U \times \GL_n(K) \hookrightarrow U \times M_n(K)$} Sei $\hat{f}$ eine Fortsetzung von $f$. Setze 
	$U = {s_{\hat{f}}}\mkern-0.5mu^{-1} \enbrace*{\Iso(E,E')}$ für $s_{\hat{f}} \colon X \to \Hom(E,E')$. Dann ist $U$ offen und $s_{\hat{f}}|_U$ entspricht einem 
	Isomorphismus $\hat{f} \colon E|_U \to E'|_U$. \bewende
\end{enumerate}
% subsection 27 (end)

\subsection[{Lemma: Trivialisierbarkeit von $E \to X \times [a,b]$, wenn $E|_{X \times [a,c]}$, $E|_{X \times [c,b]}$ trivialisierbar}]{Lemma} % (fold)
\label{sub:28}
Seien $a,b \in \mathds{R}$ mit $a <b$ und $X$ ein topologischer Raum. Sei weiter $p \colon E \to X \times [a,b]$ ein $K$-Vektorbündel. Es geben $c \in (a,b)$, so dass 
$E|_{X \times [a,c]}$ und $E|_{X \times [c,b]}$ trivialisierbar sind. Dann ist auch $E$ trivialisierbar.
\minisec{Beweis}
Sei ohne Beschränkung der Allgemeinheit $X$ zusammenhängend. Seien 
\begin{align*}
	\varphi_a \colon X \times [a,c] \times K^n &\xrightarrow{\minwidthbox{\cong}{2em}} E\big|_{X \times [a,c]} \\ 
	\varphi_b \colon X \times [c,b] \times K^n &\xrightarrow{\minwidthbox{\cong}{2em}} E\big|_{X \times [c,b]}
\end{align*}
entsprechende Trivialisierungen. Sei außerdem
\[
	h := \enbrace*{\varphi_b\big|_{X \times \set{c} \times K^n }} ^{-1} \circ  \enbrace*{\varphi_a\big|_{X \times \set{c} \times K^n}} \colon X \times \set{c} \times K^n    
\]
$h$ hat die Form $h(x,v)= (x , g(x) \cdot v)$ für eine stetige Abbildung $g \colon X \to \GL_n(K)$. Sei jetzt 
\[
	w \colon X \times [c,b] \times K^n \longrightarrow X \times [c,b] \times K^n \quad , \quad (x,t,v') \longmapsto (x,t,g(x)v')
\]
Dann gilt $\varphi_a|_{X \times \set{c} \times K^n} = \varphi_b \circ w |_{X \times \set{c} \times K^n }$. Daher können wir diese beiden Abbildungen zu einem Isomorphismus 
von Vektorbündeln $X \times [a,b] \times K^n \to E$ zusammensetzen. \bewende 
% subsection 28 (end)

\subsection[Lemma: Endliche Überdeckung $\set{U_i}_{i=1}^N$, sodass Trivialisierbarkeit auf $U_i$ gilt]{Lemma} % (fold)
\label{sub:29}
Sei $X$ ein kompakter Raum. Sei $p \colon E \to X \times I$ ein $K$-Vektorbündel. Dann gibt es eine endliche offene Überdeckung $\set{U_i}_{i=1}^N$ von $X$, so dass 
$E|_{U_i \times I}$ trivialisierbar ist.
\minisec{Beweis}
Sei $x \in X$ und $t \in I$. Dann existieren offene Umgebungen $U(t) \subset X$ und $I(t) \subset I$, sodass $E|_{U(t) \times I(t)}$ trivialisierbar ist. Da $I$ kompakt ist,
existiert eine endliche Folge $0=t_0 < t_1 < \ldots < t_N =1$ zusammen mit offenen Umgebungen $U_i$ von $x$, so dass $E|_{U_i \times [t_{i-1}, t_i]}$ trivialisierbar ist. 
Sei nun $U := \bigcap_{i=1}^N U_i$. Nach Lemma \ref{sub:28} ist $E|_{U \times I}$ trivialisierbar. Da $X \times I$ kompakt ist, wird es von endlichen vielen solcher Mengen
überdeckt. \bewende
% subsection 29 (end)

\subsection[Theorem: Isomorphismus von Vektorbündeln $E$ und $E|_{X \times \set{1}} \times I$ für $X$ kompakt]{Theorem} % (fold)
\label{sub:210}
Sei $X$ ein kompakter Raum und sei $p \colon E\to X \times I$ ein $K$-Vektorbündel. Sei $E_1 = E|_{X \times \set{1}}$.
\begin{enumerate}[a)]
	\item Es gibt einen Isomorphismus von Vektorbündeln über $X \times I$
	\[
		\begin{tikzcd}
			\varphi \colon E \rar["\cong"]  & E_1 \times I 
		\end{tikzcd}
	\]
	mit $\varphi|_{p ^{-1}(X \times \set{1})} =\id_{E_1}$
	\item Sei $(X,A)$ ein Ko-Raumpaar, $h \colon X \times I \to Y$ eine stetige Abbildung und $p' \colon E' \to Y$ ein $K$-Vektorbündel mit $E \cong h^* E'$. Falls jetzt
	$h_t(a)=h_0(a)$ für alle $a \in A$ und $t \in I$, dann lässt sich der Isomorphismus aus a) so wählen, dass er den durch $h$ induzierten Isomorphismus 
	$E|_{A \times I} \to E_1|_{A} \times I$ fortsetzt.
\end{enumerate}
\minisec{Beweis}
Nach Lemma \ref{sub:29} können wir eine Überdeckung von $X \times I$ durch offene Mengen der Form $U_i \times I$ mit $i \in \set{1,\ldots ,N}$ finden, so dass 
$E|_{U_i \times I}$ trivialisierbar ist. Seien $h_i \colon U_i \times I \times K^n \xrightarrow{\cong} E|_{U_i \times I}$ Trivialisierungen. Seien 
$\eta_i \colon X \to [0,1]$  stetige Funktionen mit folgenden Eigenschaften:

\begin{enumerate*}[(i),itemjoin={,\hspace{2em}}]
	\item $\supp(\eta_i) \subset U_i$
	\item $\max_{i=1,\ldots ,N} \eta_i(x)=1$ für alle $x \in X$
\end{enumerate*}

\noindent Sei $r_i \colon X \times I \to X \times I$ gegeben durch $r_i(x,t)=(x, \max(\eta_i(x),t))$. Ein Punkt $(x,t)$ wird durch $r_i$ entlang der $I$-Achse nach $1$ verschoben. Der 
Wert der Verschiebung wird durch $\eta_i$ festgelegt.
\missingfigure{Schöne Quadrate $X \times I$}
Nun konstruieren wir einen Morphismus von $K$-Vektorbündeln $(u_i,r_i)$ über $r_i$ wie folgt: 
\[
	u_i \enbrace*{h_i(x,t,v)} = h_i \enbrace*{r_i(x,t),v}  
\]
für $(x,t,v) \in U_i \times I \times K^n$. Wir setzen $u_i=\id$ außerhalb des Bildes von $h_i$. Wegen 
$r_i|_{X \times I\setminus{U_i} \times I} =\id_{X \times I \setminus U_i \times I}$ ist dies stetig. Außerdem kommutiert
\[
	\begin{tikzcd}
		E \rar["u_i"] \dar["p"]& R \dar["p"]\\
		X\times I \rar["r_i"] & X \times I
	\end{tikzcd}
\]
und $u_i$ ist auf den Fasern linear. Sei $r :=r_N \circ r_{N-1} \circ \ldots \circ r_1$. Wegen (ii) gilt $r(x,t)=(x,1)$. Sei $u := u_N \circ \ldots \circ u_2 \circ u_1$.
Dann ist $(u,r)$ ein Morphismus von Vektorbündeln. 
Außerdem haben wir 
\[
	\begin{tikzcd}[row sep=small]
		E \rar \drar & E \\
		& E_1 \uar[hook]
	\end{tikzcd}
\]
Setze jetzt $\varphi \colon E \to E_1 \times I$, $v \mapsto \enbrace*{u(v), \pr_I \circ p(v)} $. Dann ist $(\varphi, \id_{X \times I})$ ein Isomorphismus von Vektorbündeln 
mit $\varphi|_{X \times \set{1}} =\id$. 

Für b) sei $V_j \subset Y$ eine offene Überdeckung von $Y$, sodass $E'|_{\overline{V}_j}$ trivialisierbar ist. Sei $B_j = h ^{-1}(\overline{V}_j) \subset X \times I$. Es 
gilt 
\[
	B_j \cap (A \times I) = h_1 ^{-1}(\overline{V}_j) \cap A \times I
\]
da $h$ auf $A$ unabhängig ist von $t \in I$. Nach Konstruktion ist $E|_{A_j \times I}$ trivialisierbar. Nach Lemma \ref{sub:27} b) finden wir eine offene Umgebung 
$U_j \times I$ mit $X \supset U_j \supset A_j$, so dass $E|_{U_j \times I}$ trivialisierbar ist. Wir ergänzen die Mengen $U_j \times I$ um trivialisierende offene Umgebungen 
$U_i \times I$ mit $U_i \times A = \emptyset$, so dass $\bigcup_{i\in I} U_i \times I = X \times I$. Jetzt verwenden wir die Konstruktion aus a) für diese Überdeckung. Der 
entsprechende Isomorphismus $\varphi \colon E \to E_1 \times I$ stimmt auf $A_j \times I$ mit dem von $h$ induzierten Isomorphismus überein. \bewende 
% subsection 210 (end)
% section 2 (end)
\newpage

\section{Topologische $K$-Theorie} % (fold)
\label{sec:3}
Wir haben im letzten Kapitel folgenden kontravarianten, homotopieinvarianten Funktor kennengelernt
\[
	\begin{tikzcd}[row sep=0.4em, column sep=large,cramped]
		\TOP \rar & \MONOIDE \\
		X \rar[mapsto] & \nicefrac{\VEKT_K(X)}{\sim_{\mathrm{Iso}}} \\
		(f \colon X \to Y) \rar[mapsto] & \enbrace*{f^* \colon \nicefrac{\VEKT_K(Y)}{\sim_{\mathrm{Iso}}} \to \nicefrac{\VEKT_K(X)}{\sim_{\mathrm{Iso}}}} 
	\end{tikzcd}
\]
mit der Monoid-Operation $[E]_{\Iso} + [F]_{\Iso} = [E \oplus F]_{\Iso}$ und dem Nullvektorbündel als neutralem Element.

\minisec{Grothendieck-Konstruktion}
Sei $(M,+)$ ein kommutativer Monoid. Betrachte $\Gr(M)= \nicefrac{M \times M}{\sim}$ mit der Äquivalenzrelation
\[
	(x,y) \sim (x',y') :\iff \exists c \in M \text{ mit } x+y' +c = x' +y +c
\]
Die Idee dahinter: $[x,y]$\enquote{$=$}$x-y$. Dann ist 
\[
	\begin{tikzcd}[row sep=tiny]
		+ \colon \Gr(M) \times \Gr(M) \rar & \Gr(M) \\
		\quad \enbrace*{[x,y] , [x',y']} \rar[mapsto] & \benbrace*{x+x',y+y'}  
	\end{tikzcd}
\]
wohldefiniert. Mit dieser Operation ist $\Gr(M)$ eine abelsche Gruppe (das Inverse zu $[x,y]$ ist $[y,x]$, denn $[x+y,x+y]=[0,0]$) und $M \to \Gr(M)$, $x  \mapsto [x,0]$ ist
ein Monoidhomomorphismus.
\minisec{Beispiel}
$\Gr(\mathds{N}_0,+)= (\mathds{Z},+)$.
\minisec{Universelle Eigenschaft}
Zu jedem Monoidhomomorphismus $\varphi \colon M \to A$ in eine abelsche Gruppe $A$ existiert genau ein Gruppenhomomorphismus $\psi \colon \Gr(M) \to A$, sodass das Diagramm
\[
	\begin{tikzcd}
		M \dar\rar["\varphi"] & A \\
		\Gr(M) \urar["\exists ! \psi"', dashed]
	\end{tikzcd}
\]
kommutiert. Der Beweis ist eine Übungsaufgabe. \bigskip \\
Wir bekommen jetzt eine abelsche Gruppe $K^0(X) := \Gr \enbrace*{\nicefrac{\VEKT_{\mathds{C}}(X)}{\sim_{\Iso}}}$ (vorläufige Definition).
\minisec{Frage}
Was ist $K^n(X)$ für $n\not=0$? Was ist $K^n(X,A)$?

\subsection{Definition} % (fold)
\label{sub:31}
Sei $(X,A)$ ein Ko-Raumpaar.
\begin{enumerate}[(i)]
	\item Ein \Index{Tripel} $(E_1,f,E_0)$ über $(X,A)$ besteht aus endlichdimensionalen $\mathds{C}$-Vektorbündeln $p_i \colon E_i \to X$ und einem Isomorphismus
	$f \colon E_1|_A \to E_2|_A$.
	\item Ein \Index{Morphismus von Tripeln} $(E_1,f,E_2)$ und $(E_1',f',E_2')$ besteht aus einem Paar $(g_1,g_0)$ von Vektorbündelmorphismen $g_i \colon E_i \to E_i'$,
	sodass das folgende Diagramm kommutiert
	\[
		\begin{tikzcd}[column sep=large]
			E_1|_A \rar["g_1|_A"] \dar["f"]& E_1'|_A \dar["f'"]\\
			E_0|_A \rar["g_0|_A"] & E_0'|_A
		\end{tikzcd}
	\] 
	\item Tripel der Form $(E,\id_E,E)$ heißen \bet{elementar}\index{Tripel!elementares}.
	\item $(E_1,f,E_0)$ und $(E_1',f',E_0')$ heißen \Index{stabil isomorph}, falls 
	\[
		\enbrace*{E_1 \oplus E, f \oplus \id_E, E_0 \oplus E} \cong \enbrace*{E_1' \oplus E', f' \oplus \id_{E'}, E_0' \oplus E' } 
	\]
	Stabile Isomorphie ist eine Äquivalenzrelation (Übung). 
\end{enumerate}
% subsection 31 (end)

\subsection{Definition} % (fold)
\label{sub:32}
Sei $(X,A)$ ein Ko-Raumpaar. Dann ist $K^0(X,A)$ die Menge der stabilen Isomorphieklassen von Tripeln über $(X,A)$. Für $n>0$ definieren wir 
\[
	K^{-n}(X,A) := K^0 \enbrace[\big]{(X,A) \times (I,\partial I)^n} 
\]
Übung: $K^0(X,\emptyset)$ ist isomorph zu $\Gr \enbrace*{\nicefrac{\VEKT_\mathds{C}(X)}{\sim_{\Iso}}} $
% subsection 32 (end)
\subsection*{Bemerkung}
$\benbrace*{E_1,f,E_0}_{\stiso}=\benbrace*{0,\id,0}_\stiso $ genau dann, wenn ein Vektorbündel $E \to X$ existiert, so dass sich 
$E_1  \oplus E|_A \xrightarrow{f \oplus \id} E_0 \oplus E|_A $ zu einem Isomorphismus über ganz $X$ ausdehnen lässt.
\minisec{Beweis}
Für \enquote{$\Rightarrow$} folgt aus der Definition $(E_1 \oplus E, f \oplus \id, e_0 \oplus E) \cong \enbrace*{E', \id_{E'},E'}$. $\id_{E'}$ ist ein Isomorphismus auf ganz
$X$, also lässt sich $f \oplus \id$ zu einem Isomorphismus auf ganz $X$ ausdehnen.
\[
	\begin{tikzcd}[sep=large]
		E_1 \oplus E|_A \rar["\varphi_1|_A","\cong"'] \dar["f \oplus \id"']& E'|_A \dar["\id_{E'}","\cong"'] \\
		E_0 \oplus E|_A \rar["\cong","\varphi_0|_A"'] & E'|_A
	\end{tikzcd}
\]
Für die Rückrichtung betrachten wir $\varphi \colon E_1 \oplus E \xrightarrow{\minwidthbox{\cong}{2em}} E_0 \oplus E$, sodass $\varphi|_A = f \oplus \id$. Dann folgt
\[
	\enbrace*{E_1 \oplus E, f \oplus \id, E_0 \oplus E} \cong \enbrace*{E_ß \oplus E, \id_{E_0 \oplus E} ,E_0 \oplus E}  
\]
\[
	\begin{tikzcd}
		E_1 \oplus E|_A  \dar["f \oplus \id"] \rar["\varphi|_A"]& E_0 \oplus E|_A \dar["\id"] \\
		E_0 \oplus E|_A \rar["\id"]& E_0 \oplus E|_A
	\end{tikzcd} \bewende
\]
\subsection{Lemma} % (fold)
\label{sub:33}
Sei $(X,A)$ ein Ko-Raumpaar. Seien $p \colon E \to X$ und $p' \colon E' \to X$ $K$-Vektorbündel und $f \colon E'|_A \to E|_A$ ein Isomorphismus. Sei $f' \colon E' \to E$ ein
Isomorphismus von $K$-Vektorbündeln und sei $h \colon A \times I \to \Iso(E'|_A,E|_A)$ eine stetige Abbildung, so dass $h_t$ ein Schnitt ist für alle $t \in I$ und so dass
$h_0=f'|_A$, $h_1=f$ gilt. Dann existiert ein Isomorphismus $F \colon E' \to E$ mit $F|_A=f$ und eine Homotopie $H \colon X \times I \to \Iso(E',E)$ mit $H_0=f'$, $H_1=F$ und
$H|_{A\times I}=h$.
\minisec{Beweis}
Die Abbildungen $h$ und $f'$ liefern eine stetige Abbildung 
\[
	\begin{tikzcd}
		\tilde{h} \colon A \times I \cup X \times \set{0} \rar & \Iso \enbrace*{\pr_X^* E'|_{A \times I \cup X \times \set{0} },\pr_X^* E|_{A \times I \times X \times 
		\set{0}}} 
	\end{tikzcd}
\]
Wir müssen diese Abbildung so zu einer auf $X \times I$ ausdehnen, dass wir einen Schnitt bekommen. Betrachte \marginnote{$r$ aus Kofaserungseigenschaft}
\[
	\begin{tikzcd}
		A \times I \cup X \times \set{0} \rar["\id"] \dar[hook]& A \times I \cup X \times \set{0} \rar["\iota",hook] & X \times I\\  
		X \times I \urar[dashed,"r"]
	\end{tikzcd}
\]
Sei $r_X = \pr_X \circ \iota \circ r$ und $r_I = \pr_I \circ \iota \circ r$. Sei weiter 
\[
	\begin{tikzcd}[row sep=tiny]
		K \colon (X \times I) \times I \rar &  X \times I \\
		(x,t,s) \rar[mapsto] & \enbrace*{r_X(x,st), s r_I(x,t) + (1-s)t} 
	\end{tikzcd}
\]
Dann gilt $K_0(x,t)=(x,t)$ und $K_1(x,t)= \iota \circ r(x,t)$, also ist $\iota \circ r$ homotop zu $\id_{X \times I}$. Folglich ist
\[
	\begin{tikzcd}
		E \times I \rar["\cong"] &(\iota \circ r)^*(E \times I) \cong r^* \enbrace*{\iota^*(E \times I)} \rar["\overline{r}"] & \iota^*(E \times I) \cong E|_A \times I \cup E 
		\times \set{0} 
	\end{tikzcd}
\]
Ebenso $\overline{r}' \colon E' \times I E'|_A \times I \cup E' \times \set{0}$. $(\overline{r},r)$ und $(\overline{r}',r)$ sind Morphismen von Vektorbündeln.
Wir erhalten nun das folgende Diagramm, wobei die diagonale Abbildung rechts oben von $\tilde{h}$ induziert ist und somit ein Isomorphismus ist
\[
	\begin{tikzcd}[] % row sep=2.3em, column sep=2.8em
		E' \times I \ar[dd] \ar[rr,"\overline{r}'"] \ar[dr,"\overline{H}",dashed]& & E'|_A \times I \cup E' \times \set{0} \ar[dd] \drar["\cong"'] \\
		& E \times I \ar[rr,"\overline{r}", crossing over,pos=0.3] & & E|_A \times I \cup E \times \set{0} \ar[dd] \\
		X \times I \ar[rr,"r",pos=0.3] \ar[dr,"\id"] & & A \times I \cup X \times \set{0} \ar[dr,"\id"] \\
		& X \times I \ar[from=uu,crossing over] \ar[rr,"r"']& & A \times I \cup X \times \set{0} 
	\end{tikzcd}
\]
Die Existenz von $\overline{H}$ folgt aus der universellen Eigenschaft des Pullbacks. Es gilt nun $H=\pr_E \circ \overline{H} \colon E' \times I \to E$ und wir erhalten
eine Abbildung $X \times I \to \Iso(E' \times I, E \times I)$. $F := H|_{E' \times \set{1}} \colon E' \to E$ ist nun die gesuchte Abbildung. \bewende
% subsection 3.3 (end)
% section 3 (end)


\cleardoubleoddemptypage
\appendix
\section{Anhang}

\subsection{Lokale Trivialität des Tangentialbündels auf $S^n$} % (fold)
\label{sub:lokal_trivial_Sn}
Sei $S^n \subset \mathds{R}^{n+1}$ die $n$-Sphäre und sei 
\[
	\Tmap S^n = \set[(x,v) \in S^n \times \mathds{R}^{n+1}]{\skal*{x}{v}=0} 
\]
Wir definieren $p \colon \Tmap S^n \to S^n$ durch $p(x,v)=x$. Wir zeigen, dass $p \colon \Tmap S^n \to S^n$ lokal trivial ist.
\minisec{Beweis}
Sei $x_0 \in S^n$ gegeben. Sei $U := \set[x \in S^n]{\skal*{x}{x_0}>0}$. Da $\skal*{x_0}{x_0}=1$ ist, ist $x_0 \in U$. Weiter ist $U$ offen in $S^n$, also eine offene 
Umgebung von $x_0$. Sei $H \subset \mathds{R}^{n+1}$ die Hyperebene gegeben durch $H = \set[v \in \mathds{R}^{n+1}]{\skal*{x}{v}=0}$. Da dies ein $n$-dimensionaler 
Vektorraum ist, finden wir einen Isomorphismus $\kappa \colon H \to \mathds{R}^n$. Sei $\Tmap S^n \big|_U=p^{-1}(U)$. Wegen $\skal*{x_0}{x_0}=1$ liegt $v- \skal*{x_0}{v} x_0$
in $H$ für alle $v \in \mathds{R}^{n+1}$. Also können wir definieren:
\[
	\varphi \colon \Tmap S^n\big|_{U} \to U \times \mathds{R}^n \quad , \quad (x,v) \mapsto \enbrace[\big]{x, \kappa(v - \skal*{x_0}{v} x_0)} 
\]
Außerdem definieren wir $\psi \colon U \times \mathds{R}^n \to \Tmap S^n \big|_U$ durch
\[
	\psi(x,w) = \enbrace*{x , \kappa^{-1}(w) - \frac{\skal*{x}{\kappa^{-1}(w)}}{\skal*{x}{x_0}} x_0 } 
\]
Dies ist wohldefiniert, da für $x \in U$ der Nenner stets ungleich $0$ ist und weiter
\[
	\skal*{x}{\kappa^{-1}(w) - \frac{\skal*{x}{\kappa ^{-1}(w)}}{\skal*{x}{x_0}} x_0} = \skal*{x}{\kappa^{-1}(w)} - \frac{\skal*{x}{\kappa ^{-1}(w)}}{\skal*{x}{x_0}} 
	\skal*{x}{x_0} = 0
\]
gilt. Weiter sind beide Abbildungen offensichtlich stetig. Für die Komposition gilt
\begin{align*}
	(\psi \circ \varphi)(x,v) = \enbrace*{x, v - \skal*{x_0}{v} x_0 - \frac{\skal*{x}{v- \skal*{x_0}{v} x_0}}{\skal*{x}{x_0}}x_0 }
	= \enbrace*{x, v - \skal*{x_0}{v} x_0 + \skal*{x_0}{v}x_0} = (x,v)  
\end{align*}
Dabei haben wir benutzt, dass $\skal*{x}{v}=0$ gilt. Umgekehrt gilt
\begin{align*}
	(\varphi \circ \psi)(x,w) &= \enbrace*{x , \kappa \enbrace*{\kappa^{-1}(w) - \frac{\skal*{x}{\kappa^{-1}(w)}}{\skal*{x}{x_0}}x_0 - \skal*{x_0}{\kappa^{-1}(w) - 
	\frac{\skal*{x}{\kappa^{-1}(w)}}{\skal*{x}{x_0}}x_0} x_0 } } \\
	&= \enbrace*{x , \kappa \enbrace*{\kappa^{-1}(w) - \frac{\skal*{x}{\kappa^{-1}(w)}}{\skal*{x}{x_0}}x_0 - 
	\skal*{x_0}{\kappa^{-1}(w)} x_0 + \frac{\skal*{x}{\kappa^{-1}(w)}}{\skal*{x}{x_0}} x_0 }}\\
	&= \enbrace[\Big]{x , \kappa \enbrace*{\kappa^{-1}(w) - \skal*{x_0}{\kappa^{-1}(w)} x_0} } = (x,w)
\end{align*}
da wieder $\skal*{x_0}{\kappa ^{-1}(w)}=0$, weil $\kappa ^{-1}(w) \in H$ ist. Damit ist $\varphi$ ein Homöomorphismus. Es ist außerdem klar, dass $\varphi$ eingeschränkt auf 
die Fasern eine lineare Abbildung ist. Damit ist das Tangentialbündel $\Tmap S^n$ lokal trivial. \bewende
% subsection lokal_trivial_Sn (end)

\cleardoubleoddemptypage
\pagenumbering{Alph}
\setcounter{page}{1}

\printindex
\listoffigures
\todototoc
\listoftodos[To-do's und andere Baustellen]
\end{document}